%----------------------------------------------------------------------------------------
%   PACKAGES AND OTHER DOCUMENT CONFIGURATIONS
%----------------------------------------------------------------------------------------

\documentclass[11pt]{article}
\usepackage[left=0.75in,top=0.6in,right=0.75in,bottom=0.6in]{geometry}
\usepackage{ebgaramond}
\usepackage{enumitem}
\usepackage{array}

\pagestyle{empty}
\setlength{\parindent}{0pt}

\newcommand{\name}[1]{\begin{center}\Huge\textbf{#1}\end{center}\vspace{-0.5em}}
\newcommand{\address}[1]{\begin{center}#1\end{center}\vspace{-1em}}

\newenvironment{rSection}[1]{
  \vspace{1.5em}
  \noindent\textbf{\MakeUppercase{#1}}
  \vspace{0.2em}\hrule\vspace{0.5em}
}{}

\newenvironment{rSubsection}[4]{
  \noindent\textbf{#1} \hfill #2 \\
  \textit{#3} \hfill \textit{#4}
  \begin{itemize}[leftmargin=1.5em, itemsep=-0.2em, topsep=0pt, parsep=0pt]
}{
  \end{itemize}
  \vspace{0.5em}
}

\renewenvironment{quote}{%
  \list{}{\rightmargin=0pt \leftmargin=1.5em}%
  \item\relax
}
{\endlist}

\newenvironment{eduDetails}{%
  \begin{itemize}[label={}, leftmargin=1.5em, itemsep=0.2em, parsep=0pt, topsep=0.2em]
}{%
  \end{itemize}
}

\begin{document}

%------------------------------------------------
\name{Daniel Sabey}

\address{18a Lisa Ln \\ Ithaca, New York}
\address{(303) 356-1453 \\ danielsabey@gmail.com}

%----------------------------------------------------------------------------------------
%   EDUCATION SECTION
%----------------------------------------------------------------------------------------

\begin{rSection}{Education}
    \textbf{PhD Candidate of Economics Cornell University} \hfill \textit{August 2021 - Present}
    \begin{eduDetails}
        \item Chair: Mike Lovenheim
        \item Committee Members: Sule Alan and Evan Riehl
    \end{eduDetails}

    \textbf{MA Economics Cornell University} \hfill \textit{August 2023}
    \vspace{0.5em}
    
    \textbf{BS Economics Brigham Young University (3.97 GPA)} \hfill \textit{April 2020}
    \begin{eduDetails}
        \item Minors in Math, Business, and Korean (4.0 GPA)
        \item Brigham Young University Academic Full Tuition Scholarship for 4 Years
        \item Magna Cum Laude
    \end{eduDetails}
\end{rSection}

%----------------------------------------------------------------------------------------
%   PUBLICATIONS SECTION
%----------------------------------------------------------------------------------------

\begin{rSection}{Published Papers}
    ``Decreasing time to baccalaureate degree in the United States.'' With Denning, J. T., Eide, E. R., \& Mumford, K. J., (2022). \textit{Economics of Education Review, 90,} 102287.
    \begin{quote}
        \small{After increasing in the 1970s and 1980s, time to bachelor's degree has declined since the 1990s. We document this fact using data from three nationally representative surveys. We show that this pattern is occurring across school types and for all student types. Using administrative student records from 10 large universities, we confirm the finding and show that it is robust to alternative sample definitions. We discuss what might explain the decline in time to bachelor's degree by considering trends in student preparation, state funding, student enrollment, study time, and student employment during college.}
    \end{quote}
\end{rSection}

\begin{rSection}{Working Papers}
    ``Better Predictions, Worse Outcomes: An Experimental Evaluation of Giving Teachers Machine-Learning Forecasts of Student Achievement" With Jake Meyer.
    \begin{quote}
        \small{Teachers allocate attention and effort across students based on their own noisy assessments of who needs help. We test whether giving teachers predictions of their students' end-of-year test scores improves student achievement. In a randomized experiment across six campuses of a charter school network, treated teachers received machine-learning forecasts. The intervention successfully improved teacher prediction accuracy and increased monitoring of struggling students. However, student achievement declined by 0.099 standard deviations on average, with the largest losses among the flagged struggling students. Survey evidence suggests that outsourcing the cognitive work of diagnosing student needs crowded out the deliberation necessary for effective intervention planning.}
    \end{quote}
    
    ``Association Between Recent Medicaid Expansions and Child Maltreatment Reporting and Outcomes'' With Alex Hoagland.
    \begin{quote}
        \small{We investigate the association of recent Medicaid expansion in Maine on the volume and quality of child maltreatment reporting. We find that Medicaid expansion led to increased overall reporting, particularly from medical professionals, resulting in increased separations at lower substantiation rates. Associations were stronger among Black and female children. Our results suggest recent Medicaid expansions are associated with increased maltreatment reporting and intervention at lower specificity, resulting in a greater rate of unsubstantiated separations of children from their caregivers.}
    \end{quote}

    ``Locked Down and Lighting Up: Using COVID-19 Mobility Restrictions as an Instrument for Cigarette Stock and Smoking Intensity" With Don Kenkel and Alan Mathios.
    \begin{quote}
        \small{This paper investigates the "temptation effect" – whether increased readily available cigarette stock causally increases smoking intensity. We leverage the unprecedented natural experiment of the COVID-19 pandemic using Nielsen household scanner data, Advan Patterns cell phone data, and USAfacts.org COVID-19 data. An Interrupted Time Series (ITS) model reveals COVID having a sharp and persistent increases in both cigarette stockpiles and smoking rates. Exploration of linked survey data highlights heterogeneous responses related to stress and household characteristics. To isolate the causal effect of stock, we use county-level changes in mobility as an instrumental variable (IV) for changes in cigarette stock. The IV estimates suggest that a one-pack increase in instrumented stock causally increases daily smoking by approximately one cigarette.}
    \end{quote}
\end{rSection}

\begin{rSection}{Work in Progress}
    ``The Impact of Prisoner Education on Recidivism, Labor Market Outcomes, and Public Assistance Utilization" With Valentin Bolotnyy, Mike Lovenheim, and Natalie Millar.
    \begin{quote}
        \small{This project evaluates the impact of Career and Technical Education (CTE) programs in Tennessee state prisons on post-release outcomes. We employ an instrumental variables approach leveraging differential exposure to CTE programs based on sentence length and variation in the availability of classes across prisons to instrument for the effect of CTE enrollment.}
    \end{quote}

    ``Student Group Specific Teacher Learning'' With Jason Cook.
    \begin{quote}
        \small{This study investigates whether demographic-specific teaching skills are learnable by leveraging a natural experiment in the Columbus City Schools district. In 2003, the district shifted from race-specific to race-blind magnet school lotteries, causing significant changes to the racial composition of its schools. This paper examines how this policy change affected teacher effectiveness for different student populations and across various outcomes. Using a two-way fixed-effect model, the research analyzes how the lottery-induced changes in classroom demographics impacted teacher value-added, calculated separately before and after the policy change. By comparing these effects across different types of student outcomes (e.g., academic, behavioral), the study aims to reveal whether teachers can adapt to demographic shifts and which aspects of teaching effectiveness are most malleable mid-career.}
    \end{quote}

    ``Improving Teacher Student Matching: Comparative Advantage or Machine Learning"
    \begin{quote}
        \small{This paper explores methods to optimize teacher-student pairings by moving beyond single value-added metrics. It investigates which student characteristics (race, sex, prior performance, etc.) offer the greatest potential for gains when matching teachers based on their comparative advantages. The research introduces a novel method for classifying student "types" based on their past performance with different teachers and compares the effectiveness of traditional value-added models against higher-dimensional machine learning predictions for creating optimal classroom assignments. The paper also considers the welfare implications of reassigning students to teachers, rather than teachers to classes, and proposes an RCT to test these matching strategies in a real-world setting.}
    \end{quote}
\end{rSection}

%----------------------------------------------------------------------------------------
%   RESEARCH EXPERIENCE SECTION
%----------------------------------------------------------------------------------------

\begin{rSection}{Research Experience}

\begin{rSubsection}{Cornell University}{August 2023-Present}{Graduate Research Assistant}{Ithaca, New York}
    \item Worked with four professors:
    \subitem Mike Lovenheim (Since January 2025)
    \subitem Donald Kenkel
    \subitem Alan Mathios
    \subitem Chen Qiu
\end{rSubsection}

\begin{rSubsection}{Brigham Young University}{January 2020-July 2021}{Pre-doctoral Research Assistant}{Provo, Utah}
    \item Pre-doc with Joe Price
    \item Collaborated closely on an RCT with Phil Oreopoulos
\end{rSubsection}

\begin{rSubsection}{Brigham Young University}{April 2019-May 2020}{Research Assistant}{Provo, Utah}
    \item Worked with three professors:
    \subitem Jeff Denning
    \subitem Eric Eide
    \subitem Lars Lefgren
\end{rSubsection}

\end{rSection}

%----------------------------------------------------------------------------------------
%   TEACHING EXPERIENCE SECTION
%----------------------------------------------------------------------------------------

\begin{rSection}{Teaching Experience}

\begin{rSubsection}{Cornell Prison Education Program}{January 2026-May 2026}{Sole Instructor}{Ithaca, New York}
    \item Intro to Microeconomics
\end{rSubsection}

\begin{rSubsection}{Brigham Young University}{Summer 2025}{Sole Instructor}{Provo, Utah}
    \item Econ 110: Principles of Economics
\end{rSubsection}

\begin{rSubsection}{Cornell University}{August 2022-May 2023}{Graduate Teaching Assistant}{Ithaca, New York}
    \item Worked with two professors:
    \subitem George Boyer: The History of Economic Thought and Institutions
    \subitem Stephanie Thomas: Labor Economics
\end{rSubsection}

\begin{rSubsection}{Brigham Young University}{September 2017-September 2019}{Teaching Assistant}{Provo, Utah}
    \item Worked with 6 professors:
    \subitem Val Lambson: Game Theory and Economics
    \subitem Joe Price: Principles of Economics
    \subitem Arden Pope: Principles of Economics
    \subitem Jim Kearl: Principles of Economics
    \subitem Brennan Platt: Principles of Economics
    \subitem Kristina Bishop: Principles of Economics
\end{rSubsection}

\begin{rSubsection}{Brigham Young University}{September 2017-December 2019}{Student Consultant on Teaching}{Provo, Utah}
    \item I worked with professors to improve teaching
\end{rSubsection}

\end{rSection}

%----------------------------------------------------------------------------------------
%   SERVICE SECTION
%----------------------------------------------------------------------------------------

\begin{rSection}{Service Experience}
    \textbf{Cornell University} \hfill \textit{2024-Present}\\
    Co-Chair: Cornell Heterodox Academy Campus Community \hfill Ithaca, New York
\end{rSection}

%----------------------------------------------------------------------------------------
%   SKILLS SECTION
%----------------------------------------------------------------------------------------

\begin{rSection}{Skills}
    \begin{tabular}{@{} >{\bfseries}l @{\hspace{6ex}} l @{}}
        Coding Expertise & STATA, Python, Latex \\
        Coding Experience & SQL, MATLAB, R, Git, SLURM, Bash, SASS, Computer Vision \\
        Languages & English (native), Korean (advanced)
    \end{tabular}
\end{rSection}

%----------------------------------------------------------------------------------------
%   ACHIEVEMENTS SECTION
%----------------------------------------------------------------------------------------

\begin{rSection}{Achievements/Awards}

    Founder’s Dissertation Fellowship Award \\
    Sage Fellowship\\
    Deans List (x5)\\
    National Merit Scholar\\
    BYU Racquetball Team Member\\
    Eagle Scout

\end{rSection}

\end{document}
